\documentclass[11pt]{article}

\usepackage[french]{babel}

\usepackage[T1]{fontenc}

\usepackage{amsfonts}
\usepackage{amsmath}
\usepackage{amssymb}

\usepackage{graphicx}
\usepackage{multirow}
\usepackage{multicol}
\usepackage{xcolor}
\usepackage[shortlabels]{enumitem}

\usepackage[margin=2cm]{geometry}

\usepackage{mathtools}
\usepackage{siunitx}
\usepackage{physics}

\usepackage{tikz}
\usepackage{hyperref}

\usepackage{algorithm}
\usepackage{algpseudocode}
\usetikzlibrary{quantikz2,babel}

\tikzset{qcircuit/.style={
    column sep=.9cm,
    row sep=0.7cm,
    operator/.style={draw=black!80, line width=0.5pt, fill=gray!5, minimum height=0.7cm, minimum width=0.7cm},
    phase/.style={draw=black!80, fill=black!20, circle, minimum size=4pt},
    slice/.style={font=\footnotesize\itshape, color=gray!50, dashed}
}}


\DeclarePairedDelimiter\ceil{\lceil}{\rceil}
\DeclarePairedDelimiter\floor{\lfloor}{\rfloor}


\AtBeginDocument{\RenewCommandCopy\qty\SI}

\makeatletter
   \def\vhrulefill#1{\leavevmode\leaders\hrule\@height#1\hfill \kern\z@}
\makeatother

\newcommand{\moy}[1]{\langle #1 \rangle}
\newcommand{\pDer}[2]{\frac{\partial #1}{\partial #2}}
\newcommand{\pDerd}[2]{\frac{\partial^2 #1}{\partial #2^2}}

\newcounter{questioncounter}
\newenvironment{question}[1]{
\refstepcounter{questioncounter}
\noindent\rule{0.05\textwidth}{1pt}\,\,\raisebox{-.3\ht\strutbox}{\textbf{Question \thequestioncounter \hspace{5mm} #1}}\,\,\vhrulefill{1pt}\par\vspace{0.4cm}
}{\par}

%======================================================
%en-tete
%======================================================
\begin{document}
\selectlanguage{french}
\noindent{\Large Anthony Dubeau: Duba0207} \hfill {\Large \today}\\
{\Large Anthony Wroblewski: Wroa8915}

\par\vspace{5mm}
\begin{center}
   {\Large IFT436 - Algorithmes et structures de données} \\ \vspace{5mm}
   {\LARGE\sffamily Devoir \#1} \\ \vspace{5mm}
\end{center}

%======================================================
%question 1
%======================================================
\begin{question}{Une question d’échauffement}


On pourrait résoudre le problème en triant d'abord le tableau à l'aide du tri-fusion, puis en retournant la $k$-ième plus petite valeur ainsi que son indice original. Il faudrait donc parcourir le tableau original une premiere fois pour crée un nouveau tableau $A$ contenant des paires $(T[i],i)$, où la première composante correspond à la valeur et la seconde à son indice original dans $T$. On appliquer ensuite le tri-fusion en utilisant la premiere composante de chaque paires.\\ 

Comme le tri-fusion s'exécute en temps $O(n\log n)$ et que parcourir le tableau initial s'exécute en $O(n)$, cette approche permet de résoudre le problème en temps $O(n + n\log n) \approx O(n\log n)$.\\

Il serais peut-être possible de faire mieux en utilisant une méthode inspirée de Quicksort. Au lieu de trier complètement le tableau, on effectue seulement les partitions nécessaires pour trouver la $k$-ième plus petite valeur. Dependament des pivots, au pire on aurais quelque chose comme $O(n^2)$, mais au mieux $O(n)$.\\

\end{question}


%======================================================
%question 2
%======================================================
\begin{question}{Le problème de la sous-somme}

\begin{enumerate}[label=(\alph*)]

\item Non, cet algorithme ne résout pas correctement le problème de la sous-somme.

Considérons le contre-exemple

$$
S=\{1,2,3,4,5\}
\qquad\text{et}\qquad
t=10.
$$

Il existe une solution au problème, par exemple

$$
S'=\{1,4,5\},
$$

puisque

$$
1+4+5=10.
$$

Cependant, l'algorithme choisit toujours le plus grand élément restant dans \(S\). Il choisit donc successivement \(5\), \(4\), puis \(3\). La valeur de \(q\) évolue ainsi :

$$
q=5,\qquad q=5+4=9,\qquad q=5+4+3=12.
$$

Puisque \(q\geq t\), la boucle se termine. Comme

$$
q=12\neq 10,
$$

l'algorithme retourne \texttt{null}, alors qu'une solution existe.

l'algorithme ne résout donc pas correctement le problème de la sous-somme.

\item \textbf{(1)} Oui, cet algorithme résout correctement le problème de la sous-somme.

L'algorithme examine tous les sous-ensembles \(S' \subseteq S\). Si un sous-ensemble dont la somme vaut \(t\) existe, il sera nécessairement examiné par l'algorithme, qui le retournera. Si aucun sous-ensemble possède une somme égale à \(t\), alors l'algorithme aura examiné tous les sous-ensembles sans trouver de solution et retournera \texttt{null}. L'algorithme est donc correct.

\bigskip

\textbf{(2)} Dans le pire cas, l'algorithme entre dans la boucle une fois pour chacun des sous-ensembles de \(S\).

Pour chacun des \(n\) éléments de \(S\), il existe deux possibilités lors de la construction d'un sous-ensemble : l'élément appartient au sous-ensemble ou il n'y appartient pas. 
$$
\underbrace{2\times2\times2\times ... \times 2}_{n\times\text{fois}} = 2^n
$$
Le nombre total de sous-ensembles est donc

$$
2^n.
$$

De manière équivalente,

$$
\sum_{k=0}^{n} \binom{n}{k} = 2^n.
$$

Ainsi, dans le pire cas, l'algorithme entre dans la boucle \textbf{pour chaque} $2^n$ fois.\\

\item \textbf{(1)} L'algorithme énumère seulement les sous-ensembles de \(S\) contenant exactement \(k\) éléments.

Le nombre de façons de choisir \(k\) éléments parmi les \(n=|S|\) éléments de \(S\) est donné par le coefficient binomial

$$
\binom{n}{k}
=
\frac{n!}{k!(n-k)!}.
$$

Dans le pire cas, l'algorithme entre dans la boucle

$$
\boxed{\binom{n}{k}}
$$

fois.

\bigskip

\textbf{(2)} En géneral cette valeur est plus petite que \(2^n\), pour \(n>0\) pour $k \in \{1,2,...,\abs{S}-1\}$.

On sait que

$$
2^n
=
\sum_{i=0}^{n}\binom{n}{i}.
$$

Le terme \(\binom{n}{k}\) ne compte que les sous-ensembles de taille exactement \(k\), tandis que \(2^n\) compte tous les sous-ensembles possibles de \(S\).

Par conséquent,

$$
\boxed{\binom{n}{k} < 2^n}
$$

pour \(n>0\) et $k<n$.\\ 

Or si $k=n$ alors $\binom{n}{k}=2^n$ dans ce cas, la valeur est la même. Puisqu'on considère le pire cas, les deux algorithme respect $O(2^n)$.\\ 


Détail de notation : Je note \(\binom{n}{k}\), ou \(k\) parmi \(n\) est noter avec \(n\) en haut et \(k\) en bas. Je mentionne puisque je sais que certain font l'inverse. 

\end{enumerate}



\end{question}
%======================================================
%question 3
%======================================================
\begin{question}{Puissances en temps logarithmique}
    \begin{enumerate}[label=(\alph*)]

\item On veut démontrer que, pour tout entier \(k \geq 0\), l'algorithme $\texttt{pow(x,k)}$ retourne bien \(x^k\).

On procéde par induction sur \(k\).\\

\textbf{Cas de base.}\\

Si \(k=0\), l'algorithme retourne \(1\). Et, $x^0=1.$\\

Si \(k=1\), l'algorithme retourne \(x\). Et, $x^1=x.$\\

Donc l'algorithme est correct pour \(k=0\) et \(k=1\).

\bigskip

\textbf{Hypothèse d'induction.}

Supposons maintenant que, pour tout entier \(j<k\), l'appel

$$
\texttt{pow}(x,j)
$$

retourne correctement \(x^j\).

Considérons un entier \(k\geq 2\). L'algorithme calcule

$$
\texttt{tmp}=\texttt{pow}\left(x,\left\lfloor\frac{k}{2}\right\rfloor\right).
$$

Comme

$$
\left\lfloor\frac{k}{2}\right\rfloor < k,
$$

l'hypothèse d'induction implique que

$$
\texttt{tmp}
=
x^{\left\lfloor k/2 \right\rfloor}.
$$

Il reste alors à considérer deux cas.

\textbf{Cas 1 : \(k\) est pair.}

Il existe un entier \(m\) tel que

$$
k=2m.
$$

Dans ce cas,

$$
\left\lfloor\frac{k}{2}\right\rfloor=m,
$$

donc

$$
\texttt{tmp}=x^m.
$$

L'algorithme retourne

$$
\texttt{tmp}\cdot\texttt{tmp}
=
x^m x^m
=
x^{2m}
=
x^k.
$$

\textbf{Cas 2 : \(k\) est impair.}

Il existe un entier \(m\) tel que

$$
k=2m+1.
$$

Dans ce cas,

$$
\left\lfloor\frac{k}{2}\right\rfloor=m,
$$

donc

$$
\texttt{tmp}=x^m.
$$

L'algorithme retourne alors

$$
\texttt{tmp}\cdot\texttt{tmp}\cdot x
=
x^m x^m x
=
x^{2m+1}
=
x^k.
$$

Dans les deux cas, l'algorithme retourne correctement \(x^k\).

Par induction, l'algorithme \texttt{pow(x,k)} calcule correctement \(x^k\) pour tout \(k\geq 0\).\\


\item  Soit \(r(k)\) le nombre total d'appels à la fonction \texttt{pow}, en comptant l'appel initial.

Lorsque \(k=1\), la fonction retourne directement \(x\), donc

$$
r(1)=1.
$$

Lorsque \(k>1\), la fonction effectue exactement un appel récursif avec l'argument
\(\lfloor k/2 \rfloor\). Puisque nous supposons que \(k\) est une puissance de \(2\), \(k\) est pair et nous avons\\

\begin{equation*} 
\left\lfloor \frac{k}{2} \right\rfloor=\frac{k}{2}. \quad et \quad r(k)=1+r\left(\frac{k}{2}\right).
\end{equation*}

Nous voulons démontrer que, pour toute puissance de \(2\),

$$
r(k)=\log_2(k)+1.
$$ 
avec $k=2^m$, et \(m\geq 0\).

Nous procédons par induction sur \(m\).

\textbf{Cas de base : \(m=0\).}

On a

$$
k=2^0=1.
$$

Alors

$$
r(1)=1.
$$

Et

$$
\log_2(1)+1=0+1=1.
$$

La propriété est donc vraie pour \(m=0\).

\bigskip

\textbf{Hypothèse d'induction.}

Supposons que, pour un certain \(m \geq 0\),

$$
r(2^m)=m+1.
$$

\textbf{Étape d'induction.}

Considérons la puissance de \(2\) suivante, soit

$$
k=2^{m+1}.
$$

Pour \(k>1\), l'appel courant à \texttt{pow(x,k)} compte pour un appel et la fonction effectue ensuite un appel récursif avec \(k/2\). Par définition de \(r\), on a donc

$$
r(k)=1+r\left(\frac{k}{2}\right).
$$

En remplaçant \(k\) par \(2^{m+1}\), on obtient

$$
r(2^{m+1})
=
1+r\left(\frac{2^{m+1}}{2}\right).
$$

$$
r(2^{m+1})
=
1+r(2^m).
$$

En utilisant l'hypothèse d'induction \(r(2^m)=m+1\),

$$
r(2^{m+1})
=
1+(m+1)
=
m+2.
$$

Finalement,

$$
m+2
=
\log_2(2^{m+1})+1.
$$

% Ainsi,

% $$
% r(2^{m+1})
% =
% \log_2(2^{m+1})+1.
% $$


% Par l'hypothèse d'induction,

% $$
% r(2^{m+1})
% =
% 1+(m+1)
% =
% m+2.
% $$

% Or,

% $$
% \log_2(2^{m+1})+1
% =
% m+1+1
% =
% m+2.
% $$

On peut remplacer,

$$
r(2^{m+1})
=
\log_2(2^{m+1})+1.
$$

Par induction, pour toute puissance de \(2\),

$$
\boxed{r(k)=\log_2(k)+1}.
$$

\end{enumerate}
    
\end{question}

\bibliographystyle{apsrev4-2}
\bibliography{references}
\end{document}
